%==============================================================================
%  Processeur RISC-V (RV32I) pipeliné -- Rapport technique
%==============================================================================
\documentclass[11pt,a4paper,oneside]{report}

% ---- Geometrie / espacement -------------------------------------------------
\usepackage[a4paper,left=2.8cm,right=2.8cm,top=2.6cm,bottom=2.8cm]{geometry}
\usepackage{setspace}
\onehalfspacing

% ---- Encodage / polices -----------------------------------------------------
\usepackage[utf8]{inputenc}
\usepackage[T1]{fontenc}
\usepackage{lmodern}
\usepackage[french]{babel}
\usepackage{microtype}

% ---- Maths / graphiques / tableaux ------------------------------------------
\usepackage{amsmath,amssymb}
\usepackage{graphicx}
\usepackage{booktabs}
\usepackage{tabularx}
\usepackage{array}
\usepackage{multirow}
\usepackage{float}
\usepackage{siunitx}

% ---- Mise en forme des chapitres / sections ---------------------------------
\usepackage{titlesec}
\titleformat{\chapter}[display]
  {\normalfont\Large\bfseries}{Chapitre \thechapter}{0.6em}{\Huge}
\titlespacing*{\chapter}{0pt}{-30pt}{20pt}

% ---- Listings (Verilog) avec coloration syntaxique avancée ------------------
\usepackage{listings}
\usepackage{xcolor}

% Couleurs personnalisées pour la coloration syntaxique
\definecolor{codebackground}{RGB}{248,248,248}
\definecolor{codekeyword}{RGB}{0,0,180}
\definecolor{codecomment}{RGB}{0,128,0}
\definecolor{codestring}{RGB}{163,21,21}
\definecolor{codenumber}{RGB}{128,128,128}
\definecolor{codeidentifier}{RGB}{0,0,0}
\definecolor{codedirective}{RGB}{128,0,128}

\lstdefinelanguage{Verilog}{
  morekeywords=[1]{module,endmodule,input,output,inout,wire,reg,assign,always,
                posedge,negedge,begin,end,if,else,case,casez,casex,endcase,
                parameter,localparam,default,for,generate,endgenerate,
                function,endfunction,task,endtask,initial,integer,
                signed,unsigned},
  morekeywords=[2]{`define,`ifdef,`ifndef,`else,`endif,`include,`timescale,
                `undef,`default\_nettype},
  sensitive=true,
  morecomment=[l]{//},
  morecomment=[s]{/*}{*/},
  morestring=[b]"
}
\lstset{
  language=Verilog,
  basicstyle=\ttfamily\footnotesize,
  keywordstyle=[1]\bfseries\color{codekeyword},
  keywordstyle=[2]\color{codedirective},
  commentstyle=\itshape\color{codecomment},
  stringstyle=\color{codestring},
  identifierstyle=\color{codeidentifier},
  numbers=left, numberstyle=\tiny\color{codenumber}, numbersep=8pt,
  backgroundcolor=\color{codebackground},
  frame=single, framerule=0.4pt, rulecolor=\color{gray!50},
  breaklines=true, columns=fullflexible,
  captionpos=b, xleftmargin=2em, framexleftmargin=1.5em,
  showstringspaces=false,
  tabsize=4,
  aboveskip=1em, belowskip=1em
}

% ---- Hyperliens / references ------------------------------------------------
\usepackage{hyperref}
\hypersetup{
  colorlinks=true, linkcolor=black, citecolor=black, urlcolor=blue!50!black,
  pdftitle={Processeur RISC-V RV32I pipeline -- Rapport technique},
  pdfauthor={Nom de l'auteur}
}
% cleveref retiré -- on utilise \ref{} standard pour éviter tout conflit


%==============================================================================
\begin{document}

% ---- Pages liminaires (numerotation romaine) --------------------------------
\pagenumbering{roman}

\begin{titlepage}
  \centering
  \vspace*{2.5cm}
  {\Large\bfseries Conception et vérification d'un\\[0.3em]
   processeur RISC-V RV32I pipeliné en Verilog\par}
  \vspace{2cm}
  {\large Rapport technique\par}
  \vspace{2cm}
  {\large Nom de l'auteur\par}
  {\normalsize Département de génie électrique et informatique\par}
  {\normalsize Université / Institution\par}
  \vspace{2cm}
  {\normalsize \today\par}
  \vfill
\end{titlepage}

\chapter*{Résumé}
\addcontentsline{toc}{chapter}{Résumé}
Ce rapport présente la conception, l'implémentation et la
vérification fonctionnelle d'un c\oe{}ur de processeur RISC-V pipeliné
à cinq étages mettant en \oe{}uvre le jeu d'instructions de base
RV32I. La microarchitecture suit l'organisation canonique
Récupération d'instruction (IF), Décodage (ID), Exécution (EX),
Accès mémoire (MEM) et Écriture en retour (WB), augmentée d'une
unité de renvoi des opérandes, d'un détecteur de blocage pour les
dépendances charge-utilisation et d'un mécanisme de vidage du
pipeline destiné à atténuer les aléas de données et de contrôle.
La conception est décrite en Verilog-2001 synthétisable, organisée
en modules par étage clairement délimités sous un module
enveloppe unique de plus haut niveau. La vérification
fonctionnelle est réalisée à l'aide de la suite officielle
\texttt{riscv-tests} \texttt{rv32ui-p-*}, chaque test étant compilé
en image hexadécimale puis rejoué sur un banc de test
auto-vérifiant. Le rapport documente l'architecture, la stratégie
de gestion des aléas, les interfaces Verilog de chaque module et
les résultats de la vérification.

\tableofcontents
\listoffigures
\listoftables

% ---- Corps du rapport (numerotation arabe) ----------------------------------
\clearpage
\pagenumbering{arabic}


%==============================================================================
\chapter{Introduction}
\label{ch:intro}

\section{Motivation et contexte}
Le jeu d'instructions RISC-V (Reduced Instruction Set Computer~V)
s'est imposé au cours de la dernière décennie comme un standard
ouvert et libre pour la conception de processeurs, offrant une
alternative claire, modulaire et extensible aux jeux d'instructions
propriétaires qui dominent depuis longtemps les marchés de
l'embarqué et du calcul généraliste. Son sous-ensemble entier de
base, RV32I, est volontairement réduit---moins de cinquante
instructions---tout en restant suffisant pour supporter un modèle
de programmation généraliste complet ainsi que les chaînes de
compilation modernes, les noyaux de système d'exploitation et les
suites de tests de référence. La conjonction de l'ouverture, de la
simplicité et de la maturité de la chaîne d'outils a fait de RV32I
un véhicule particulièrement adapté à l'étude de la conception
classique de processeurs pipelinés, sans les contraintes
d'encodages historiques ni les restrictions de licence.

Dans un cadre académique, la construction d'un c\oe{}ur RV32I
fonctionnel rassemble des sujets habituellement enseignés
isolément : encodage et décodage du jeu d'instructions, conception
de machines à états finis, description au niveau transfert de
registres (RTL), analyse des aléas et vérification fonctionnelle.
Le présent travail aborde l'ensemble de ces objectifs au sein
d'un même artefact intégré : une description Verilog synthétisable
d'un c\oe{}ur RV32I pipeliné à cinq étages, accompagnée d'un banc de
test auto-vérifiant et d'une suite de non-régression issue du
projet officiel \texttt{riscv-tests}.

\section{Objectifs du projet}
Le travail documenté dans ce rapport poursuit trois objectifs
principaux :
\begin{itemize}
  \item concevoir un c\oe{}ur synthétisable, en ordre, à émission
        unique, pipeliné à cinq étages, exécutant correctement
        l'intégralité du jeu d'instructions utilisateur RV32I ;
  \item implémenter un sous-système complet de gestion des aléas
        comprenant le renvoi (\textit{forwarding}) des opérandes,
        l'insertion de blocages pour les dépendances de type
        charge-utilisation, ainsi que le vidage du pipeline lors
        des branchements et des sauts pris ;
  \item vérifier la conception résultante à l'aide de la suite de
        non-régression officielle \texttt{riscv-tests}
        \texttt{rv32ui-p-*}, fournissant ainsi une preuve
        objective et reproductible de la conformité au jeu
        d'instructions.
\end{itemize}
Un objectif secondaire, d'ordre pédagogique, est la clarté de la
lecture : le code source est délibérément découpé en un module
Verilog par étage du pipeline, de sorte que la conception puisse
être abordée de haut en bas, chaque étage pouvant être compris
isolément avant l'examen de la logique de gestion des aléas qui
relie les étages entre eux.


\section{Périmètre et limitations}
L'implémentation actuelle ne cible que le sous-ensemble utilisateur
de RV32I. Les registres de contrôle et de statut (CSR),
l'architecture privilégiée, la gestion des trappes et des
interruptions, ainsi que les extensions M (multiplication et
division), A (atomique), F/D (virgule flottante) et C (compressée)
sont explicitement hors du champ de l'étude. Les instructions
exceptionnelles (\texttt{ECALL}, \texttt{EBREAK}) et les
écritures mémoire mal alignées sont détectées et signalées par un
signal \texttt{exception\_o} de plus haut niveau, mais ne sont pas
traitées par un gestionnaire de trappe interne au c\oe{}ur ; leur
prise en charge incombe au système sur puce (SoC) environnant. De
même, le sous-système mémoire est modélisé sous forme de ports
idéaux IMEM et DMEM de type Harvard ; les caches, les protocoles
de bus (par exemple AXI) et les contrôleurs mémoire ne sont pas
abordés.

\section{Contributions}
Les contributions principales de ce travail sont :
\begin{enumerate}
  \item une implémentation Verilog modulaire, claire et
        entièrement synthétisable d'un pipeline RV32I à cinq
        étages ;
  \item un sous-système documenté de gestion des aléas combinant
        des sélecteurs de renvoi sur deux bits (\texttt{forwardA},
        \texttt{forwardB}), un sélecteur de renvoi MEM-vers-MEM
        sur un bit pour la donnée à écrire (\texttt{forwardM}),
        un blocage charge-utilisation (\texttt{stall\_IF\_ID})
        ainsi que des vidages du pipeline (\texttt{flush\_IF\_ID},
        \texttt{flush\_EX}) ;
  \item un flot de vérification reproductible bâti autour de la
        suite officielle \texttt{riscv-tests}
        \texttt{rv32ui-p-*}, accompagné d'un décompte
        succès/échec par test.
\end{enumerate}

\section{Organisation du rapport}
Le \cref{ch:arch} introduit l'architecture de plus haut niveau et
la hiérarchie des modules de la présente implémentation. Le
\cref{ch:pipeline} détaille chacun des cinq étages du pipeline.
Le \cref{ch:hazards} décrit la logique de gestion des aléas. Le
\cref{ch:impl} discute des conventions d'implémentation Verilog,
et le \cref{ch:verif} expose la méthodologie et les résultats de
la vérification. Le \cref{ch:discussion} aborde les limitations
et les optimisations envisageables, et le \cref{ch:conclusion}
conclut le rapport.


%==============================================================================
\chapter{Aperçu de l'architecture}
\label{ch:arch}

\section{Organisation de plus haut niveau}
Le processeur est encapsulé dans un unique module de plus haut
niveau, \texttt{mpt\_top}, qui expose une horloge, une
réinitialisation active à l'état bas, une entrée d'adresse de
démarrage ainsi que deux interfaces mémoire de type Harvard : un
port en lecture seule vers la mémoire d'instructions (IMEM) et un
port en lecture/écriture, à granularité d'octet, vers la mémoire
de données (DMEM). Une sortie unique \texttt{exception\_o} signale
l'occurrence d'un \texttt{ECALL}, d'un \texttt{EBREAK} ou d'une
écriture mémoire mal alignée. L'interface du plus haut niveau est
intentionnellement minimale, de manière à ce que le c\oe{}ur puisse
être intégré dans un système sur puce (SoC) plus large doté d'un
sous-système mémoire arbitraire, qu'il s'agisse de mémoire bloc
sur puce, d'un contrôleur DRAM externe placé derrière un cache
ou d'un pont vers une matrice de périphériques.

\begin{figure}[H]
  \centering
  \rule{\textwidth}{0.4pt}\\[2cm]
  \textit{[Insérer image ici]}\\[2cm]
  \rule{\textwidth}{0.4pt}
  \caption{Schéma bloc de plus haut niveau du processeur montrant les
  entrées-sorties de \texttt{mpt\_top}, les interfaces IMEM et DMEM,
  l'entrée d'adresse de démarrage et la sortie \texttt{exception\_o}.}
  \label{fig:top-level}
\end{figure}

\section{Hiérarchie des modules}
En interne, \texttt{mpt\_top} instancie un module par étage du
pipeline---\texttt{mpt\_IF\_stage}, \texttt{mpt\_ID\_stage},
\texttt{mpt\_EX\_stage}, \texttt{mpt\_MEM\_stage} et
\texttt{mpt\_WB\_stage}---ainsi qu'un petit nombre de blocs
combinatoires partagés. \texttt{mpt\_decoder} et
\texttt{mpt\_control} produisent conjointement le mot de contrôle
propre à chaque instruction à partir du mot de 32 bits.
\texttt{mpt\_regfile} implémente le banc de registres
architectural de 32~$\times$~32 bits, à lecture asynchrone et
écriture synchrone, le registre \texttt{x0} étant câblé à zéro.
\texttt{mpt\_alu} réalise l'ensemble des opérations
arithmétiques, logiques et de décalage. \texttt{mpt\_branchgen}
résout la direction des branchements et calcule les adresses
cibles des branchements et des sauts. Les encodages du jeu
d'instructions, les codes opérations de l'ALU et les largeurs des
paramètres sont centralisés dans le fichier d'en-tête
\texttt{mpt\_definitions.vh}, inclus par chaque module afin que
toute modification d'un opcode ou d'une largeur se propage de
façon cohérente à l'ensemble de la conception.

\begin{figure}[H]
  \centering
  \rule{0.85\textwidth}{0.4pt}\\[2cm]
  \textit{[Insérer image ici]}\\[2cm]
  \rule{0.85\textwidth}{0.4pt}
  \caption{Hiérarchie des modules de \texttt{mpt\_top}, listant les cinq
  modules d'étage et les blocs partagés de décodage, de contrôle, de banc
  de registres, d'ALU et de génération de branchements.}
  \label{fig:hierarchy}
\end{figure}


\section{Aperçu du chemin de données}
Le chemin de données suit l'organisation classique en cinq
étages. À l'étage IF, le compteur ordinal (PC) est présenté à
l'IMEM et le mot retourné est verrouillé dans le registre
pipeline IF/ID. À l'étage ID, l'instruction est décodée,
l'immédiat est étendu par signe ou par zéro selon son format
(I/S/B/U/J), et les deux opérandes source sont lus dans le banc
de registres. À l'étage EX, l'ALU combine ses opérandes---
choisis parmi les sorties du banc de registres, l'immédiat et les
valeurs renvoyées---pour calculer soit un résultat arithmétique,
soit une adresse mémoire effective ; en parallèle,
\texttt{mpt\_branchgen} évalue la condition de branchement et
calcule l'adresse cible. À l'étage MEM, les chargements et les
écritures sont émis vers la DMEM avec le masque d'octets
approprié ; les chargements sous-mot (\texttt{LB}, \texttt{LH},
\texttt{LBU}, \texttt{LHU}) sont étendus par signe ou par zéro
sur la largeur complète d'un registre. Enfin, à l'étage WB, le
résultat---sortie de l'ALU, donnée chargée ou adresse de retour
$\texttt{PC}+4$---est écrit en retour dans le banc de registres.

\begin{figure}[H]
  \centering
  \rule{\textwidth}{0.4pt}\\[2cm]
  \textit{[Insérer image ici]}\\[2cm]
  \rule{\textwidth}{0.4pt}
  \caption{Chemin de données pipeliné du c\oe{}ur RV32I montrant les étages
  IF, ID, EX, MEM et WB, les quatre registres pipeline inter-étages ainsi
  que les multiplexeurs principaux.}
  \label{fig:datapath}
\end{figure}

\section{Aperçu du chemin de contrôle}
Les signaux de contrôle sont produits une seule fois par
instruction à l'étage ID, puis circulent aux côtés des opérandes
à travers les registres pipeline. Le décodeur
(\texttt{mpt\_decoder}) classifie l'instruction d'après son code
opération et ses champs \texttt{funct3}/\texttt{funct7}, et émet
les autorisations de lecture et d'écriture du banc de registres,
le sélecteur de format d'immédiat, le sélecteur de source de
l'ALU (\texttt{ID\_alu\_source\_sel}), le code opération de l'ALU
(\texttt{ID\_alu\_ctrl}), le type de branchement
(\texttt{ID\_branch\_op}), l'opération mémoire
(\texttt{ID\_mem\_op}) ainsi que le sélecteur de source de
l'écriture en retour (\texttt{ID\_memtoreg}). Un bloc dédié de
détection des aléas, décrit au \cref{ch:hazards}, active les
signaux \texttt{stall\_IF\_ID}, \texttt{flush\_IF\_ID},
\texttt{flush\_EX} ainsi que les sélecteurs de renvoi
\texttt{forwardA}, \texttt{forwardB} et \texttt{forwardM}.

\begin{figure}[H]
  \centering
  \rule{\textwidth}{0.4pt}\\[2cm]
  \textit{[Insérer image ici]}\\[2cm]
  \rule{\textwidth}{0.4pt}
  \caption{Chemin de contrôle montrant la génération du mot de contrôle à
  l'étage ID puis sa propagation à travers les registres pipeline ID/EX,
  EX/MEM et MEM/WB.}
  \label{fig:controlpath}
\end{figure}

\section{Interfaces mémoire}
Les deux ports mémoire sont adressés par mot et larges de 32
bits. Le port IMEM fournit une lecture combinatoire du mot
d'instruction à l'adresse \texttt{IMEM\_addr\_o} ; l'adresse de
démarrage est fournie de l'extérieur par \texttt{boot\_addr\_i}
et chargée dans le PC à la réinitialisation. Le port DMEM gère
les écritures à granularité d'octet, de demi-mot et de mot grâce
au signal de validation sur quatre bits
\texttt{DMEM\_wr\_byte\_en\_o} ; les lectures retournent toujours
un mot de 32 bits dont l'étage MEM extrait et étend la portion
sous-mot demandée. Les écritures mémoire mal alignées sont
détectées à l'étage MEM et signalées sur \texttt{exception\_o}.

\section{Paramétrage}
Le c\oe{}ur est paramétré au moyen de \texttt{mpt\_definitions.vh}
et de la liste de paramètres du module de plus haut niveau,
exposant la largeur des données de registre
(\texttt{REG\_DATA\_WIDTH}, par défaut 32), la largeur d'adresse
du banc de registres (\texttt{REGFILE\_ADDR\_WIDTH}, par défaut
5), la largeur d'adresse de l'IMEM (\texttt{IMEM\_ADDR\_WIDTH},
par défaut 32) ainsi que la largeur du code opération de l'ALU
(\texttt{ALU\_OP\_WIDTH}, par défaut 4). Bien que le présent
travail vise RV32I, ce paramétrage laisse ouverte la possibilité
d'une migration ultérieure vers RV64I sans refonte intrusive du
code.


%==============================================================================
\chapter{Conception des étages du pipeline}
\label{ch:pipeline}

Ce chapitre détaille successivement les cinq étages du pipeline,
puis les trois blocs combinatoires partagés par plusieurs étages :
l'unité arithmétique et logique (\texttt{mpt\_alu}), le banc de
registres (\texttt{mpt\_regfile}) et le générateur d'adresses de
branchement (\texttt{mpt\_branchgen}). Chaque étage est présenté
suivant une trame uniforme : rôle, interfaces, registres pipeline
en sortie, comportement combinatoire interne et points d'attention
microarchitecturaux propres à l'implémentation.

\section{Étage IF -- Récupération d'instruction}
\label{sec:if}

\subsection{Rôle et interfaces}
Le module \texttt{mpt\_IF\_stage} héberge le compteur ordinal
\texttt{IMEM\_addr\_o}, qui sert simultanément d'adresse présentée
à la mémoire d'instructions et de registre architectural du PC. À
chaque front montant de \texttt{clk\_i}, le module met à jour
\texttt{IMEM\_addr\_o}, puis recopie son ancienne valeur dans
\texttt{IF\_pc\_o}. Cette duplication apparente---deux registres
représentant le PC---est en réalité essentielle au pipeline : elle
permet de propager vers ID l'adresse de l'instruction réellement
fetchée pendant que IF est déjà en train de calculer l'adresse de
l'instruction suivante.

\subsection{Multiplexeur d'entrée du PC}
Le c\oe{}ur de l'étage est un multiplexeur 4-vers-1
combinatoire qui sélectionne la valeur de \texttt{pc\_next} :
\begin{enumerate}
  \item priorité~1 : \texttt{ID\_jump\_en\_i} actif --- saut
        inconditionnel (\texttt{JAL} ou \texttt{JALR}) résolu en
        ID, le PC suivant est l'adresse cible
        \texttt{BG\_pc\_dest\_i} fournie par le générateur de
        branchement ;
  \item priorité~2 : \texttt{EX\_branch\_en\_i} actif ---
        branchement conditionnel pris confirmé en EX, le PC
        suivant est \texttt{EX\_pc\_dest\_i} ;
  \item priorité~3 : \texttt{stall\_i} actif --- aléa
        charge-utilisation, on doit recharger l'instruction
        courante, donc \texttt{pc\_next = IMEM\_addr\_o - 4} ;
  \item cas par défaut : avancement séquentiel
        \texttt{pc\_next = IMEM\_addr\_o + 4}.
\end{enumerate}
La priorité~3 mérite une explication. Comme \texttt{IMEM\_addr\_o}
a déjà été incrémenté lors du cycle où l'aléa est détecté, le
processeur a en fait dépassé l'instruction qui doit être rejouée.
La soustraction de quatre rétablit cette adresse au cycle suivant,
ce qui équivaut fonctionnellement à figer le PC le temps qu'une
bulle traverse l'étage ID.

\subsection{Injection de NOP par le signal de \textit{flush}}
La sortie \texttt{IF\_instruction\_o} n'est pas un registre
synchrone supplémentaire mais un multiplexeur 2-vers-1
combinatoire piloté par \texttt{flush\_i} :
\[
\texttt{IF\_instruction\_o} =
\begin{cases}
\texttt{32'h0000\_0000} & \text{si } \texttt{flush\_i} = 1 \\
\texttt{IMEM\_data\_i} & \text{sinon.}
\end{cases}
\]
La constante \texttt{32'h00000000} se décode en
\texttt{ADDI x0, x0, 0}, soit l'encodage canonique du NOP RISC-V.
Cette technique, courante mais subtile, permet de transformer
n'importe quelle instruction fetchée en un NOP architectural sans
ajouter de chemin de contrôle dédié à l'étage ID.

\subsection{Comportement à la réinitialisation}
Au front montant suivant l'assertion de \texttt{resetn\_i} bas,
\texttt{IMEM\_addr\_o} est forcé à \texttt{boot\_addr\_i} et
\texttt{IF\_pc\_o} est mis à zéro. La quasi-totalité des suites
\texttt{rv32ui-p} sont liées à l'adresse de démarrage
\texttt{0x80000000} mais sont relogées à \texttt{0x00000000} dans
notre flot, de sorte que le testbench fournit
\texttt{boot\_addr\_i = 32'b0}.

\begin{figure}[H]
  \centering
  \rule{0.9\textwidth}{0.4pt}\\[2cm]
  \textit{[Insérer image ici]}\\[2cm]
  \rule{0.9\textwidth}{0.4pt}
  \caption{Détail de l'étage IF montrant le multiplexeur 4-vers-1
  d'entrée du PC, la duplication \texttt{IMEM\_addr\_o}/\texttt{IF\_pc\_o}
  et le multiplexeur d'injection de NOP piloté par \texttt{flush\_i}.}
  \label{fig:if-stage}
\end{figure}


\section{Étage ID -- Décodage et lecture des registres}
\label{sec:id}

\subsection{Architecture en enveloppe}
Le module \texttt{mpt\_ID\_stage} est une enveloppe structurelle
combinant deux sous-blocs entièrement combinatoires
(\texttt{mpt\_decoder} et le port de lecture de
\texttt{mpt\_regfile}) avec les flip-flops du registre pipeline
ID/EX. Cette organisation présente un avantage de lisibilité :
toute la logique de décodage, sans état, vit dans un fichier
indépendant ; tout l'état architectural transite par les
toujours-blocs synchrones de l'enveloppe.

\subsection{Sous-bloc 1 -- Décodeur}
\texttt{mpt\_decoder} reçoit le mot de 32 bits
\texttt{IF\_instruction\_i} ainsi que \texttt{IF\_pc\_i}, et
produit en parallèle :
\begin{itemize}
  \item les trois adresses de registre \texttt{rd\_addr\_o},
        \texttt{rs1\_addr\_o}, \texttt{rs2\_addr\_o} ;
  \item les cinq versions étendues d'immédiat
        \texttt{IMM\_I}, \texttt{IMM\_S}, \texttt{IMM\_B},
        \texttt{IMM\_U} et \texttt{IMM\_J} ;
  \item le sélecteur d'origine d'opérande
        \texttt{alu\_source\_sel\_o} (deux bits indiquant
        respectivement si l'opérande~1 et l'opérande~2 doivent
        être un immédiat) ;
  \item le code opération de l'ALU \texttt{alu\_op\_o} (4 bits) ;
  \item les signaux de branchement \texttt{branch\_op\_o} et
        \texttt{branch\_flag\_o} ;
  \item les autorisations mémoire \texttt{mem\_wr\_en\_o},
        \texttt{mem\_rd\_en\_o}, le code d'opération mémoire
        \texttt{mem\_op\_o} et le sélecteur d'écriture
        \texttt{memtoreg\_o} ;
  \item l'autorisation d'écriture \texttt{rd\_wr\_en\_o} et le
        drapeau \texttt{exception\_o}.
\end{itemize}

\subsubsection*{Construction des immédiats}
La spécification RV32I retient cinq formats d'immédiat (I, S, B, U,
J) qui ne diffèrent que par la position des bits dans le mot
d'instruction. Le décodeur tire parti du fait que le bit de signe
\texttt{instruction\_i[31]} est \emph{toujours} le bit le plus
significatif de l'immédiat, quel que soit le format. Cette
particularité, voulue par les architectes de RISC-V, permet
d'implémenter l'extension de signe par une réplique unique du même
fil sur 12 ou 20 bits, indépendamment de l'opcode :
\begin{lstlisting}[caption={Construction des cinq immédiats dans
\texttt{mpt\_decoder}.},label={lst:imm}]
IMM_I = { {20{instruction_i[31]}}, instruction_i[31:20] };
IMM_S = { {20{instruction_i[31]}}, instruction_i[31:25],
          instruction_i[11:7] };
IMM_B = { {20{instruction_i[31]}}, instruction_i[7],
          instruction_i[30:25], instruction_i[11:8], 1'b0 };
IMM_U = { instruction_i[31:12], 12'b0 };
IMM_J = { {12{instruction_i[31]}}, instruction_i[19:12],
          instruction_i[20], instruction_i[30:25],
          instruction_i[24:21], 1'b0 };
\end{lstlisting}
La concaténation finale d'un \texttt{1'b0} dans \texttt{IMM\_B} et
\texttt{IMM\_J} reflète l'alignement obligatoire des cibles de
branchement et de saut sur deux octets, qui n'est pas codé dans
l'instruction.

\subsubsection*{Décodage des opcodes}
Le décodeur examine \texttt{OPCODE = instruction\_i[6:0]} dans une
unique structure \texttt{case}. Pour chacune des huit familles
d'opcode (\texttt{OP}, \texttt{OP\_IMM}, \texttt{BRANCH},
\texttt{LUI}, \texttt{AUIPC}, \texttt{JAL}, \texttt{JALR},
\texttt{LOAD}, \texttt{STORE}), il configure simultanément le code
ALU, l'origine des opérandes, l'immédiat à propager et les
autorisations de mémoire et d'écriture en retour. Trois détails
méritent d'être soulignés :
\begin{enumerate}
  \item \texttt{ADD} et \texttt{SUB} (ainsi que \texttt{SRL} et
        \texttt{SRA}) partagent le même \texttt{FUNCT3} et ne se
        distinguent que par le bit \texttt{FUNCT7 = instruction\_i[30]}
        (les autres bits du champ \texttt{funct7} sont ignorés) ;
  \item l'autorisation d'écriture \texttt{rd\_wr\_en\_o} est
        explicitement masquée à zéro lorsque
        \texttt{instruction\_i[11:7] == 0}, c'est-à-dire lorsque
        l'instruction tente d'écrire dans \texttt{x0}, ce qui
        soulage l'unité de gestion des aléas (un faux positif
        \texttt{x0} ne pourra jamais déclencher un renvoi) ;
  \item \texttt{JAL} et \texttt{JALR} forcent
        \texttt{alu\_source\_sel\_o = 2'b10}, ce qui aiguille
        l'opérande~1 vers \texttt{pc\_i}. La constante~4 est, elle,
        injectée plus tard dans l'étage EX ; ainsi le calcul de
        l'adresse de retour \texttt{PC + 4} est partagé entre le
        décodeur et l'étage EX, sans additionneur dédié.
\end{enumerate}
Enfin, \texttt{exception\_o} est levé exactement lorsque le mot
d'instruction est l'un des deux encodages \texttt{ECALL}
(\texttt{32'h00000073}) ou \texttt{EBREAK}
(\texttt{32'h00100073}).


\subsection{Sous-bloc 2 -- Banc de registres}
Le module \texttt{mpt\_regfile} expose deux ports de lecture
combinatoires (\texttt{rs1\_data\_o}, \texttt{rs2\_data\_o}) et un
port d'écriture synchrone. Sa particularité est de réaliser, sur
chaque port de lecture, un renvoi interne en lecture-vers-écriture
combinatoire : si l'adresse demandée à un port de lecture est
identique à l'adresse en cours d'écriture (\texttt{rd\_addr\_i}) et
que \texttt{rd\_wr\_en\_i} est actif, le multiplexeur substitue
combinatoirement \texttt{rd\_wr\_data\_i} à la sortie du tableau
mémoire. Cette astuce résout l'aléa structurel
\og{}écriture-lors-de-lecture\fg{} sans coût additionnel en logique
séquentielle. L'éligibilité \texttt{rd\_addr\_i != 0} garantit
qu'aucune écriture vers \texttt{x0} ne traverse jamais ce chemin
de renvoi.

\subsection{Pipeline ID/EX et injection de bulles}
Trois toujours-blocs synchrones distincts forment le registre
pipeline ID/EX :
\begin{itemize}
  \item \emph{Bloc A} : capture les adresses de registre, les
        immédiats et le drapeau d'exception. Ces signaux ne sont
        jamais masqués par \texttt{stall\_i} : ils ne déclenchent
        aucune action et peuvent être traités comme du
        \og{}pass-through inoffensif\fg{}.
  \item \emph{Bloc B} : capture les signaux de contrôle. Ce bloc
        implémente la \emph{politique d'insertion de bulle} :
        \begin{itemize}
          \item sur \texttt{flush\_i} ou
                \texttt{resetn\_i = 0}, tous les signaux sont
                effacés ;
          \item sur \texttt{stall\_i = 1}, les signaux passifs
                (PC, code ALU, immédiats) sont gelés, mais les
                trois autorisations destructives
                \texttt{ID\_mem\_wr\_en\_o},
                \texttt{ID\_mem\_rd\_en\_o},
                \texttt{ID\_rd\_wr\_en\_o} sont forcées à zéro.
                C'est cette politique qui transforme l'instruction
                rejouée en une bulle sans état pendant le cycle de
                blocage ;
          \item dans le cas standard, les sorties du décodeur sont
                propagées sans modification.
        \end{itemize}
  \item \emph{Bloc C} : recopie synchroniquement
        \texttt{rs1\_data} et \texttt{rs2\_data} dans
        \texttt{rs1\_rd\_data} et \texttt{rs2\_rd\_data}. Cette
        registration intermédiaire permet d'isoler le banc de
        registres du chemin combinatoire interne au bloc D
        ci-dessous.
\end{itemize}

\subsection{Renvoi interne ID-WB}
Le bloc combinatoire D effectue un renvoi de fin de pipeline
strictement interne à l'étage ID : si l'instruction qui finalise
sa phase WB le cycle courant écrit dans le même registre que
\texttt{rs1} ou \texttt{rs2} de l'instruction en train d'être
décodée, et que cette écriture porte sur un registre non nul,
\texttt{ID\_rs1\_data\_o} ou \texttt{ID\_rs2\_data\_o} reçoit
combinatoirement \texttt{WB\_rd\_wr\_data\_i}. Ce mécanisme évite
qu'une donnée venant d'être écrite par WB et lue le même cycle ne
soit obsolète---l'écriture étant synchrone, elle ne sera visible
au tableau mémoire qu'au cycle suivant.

\begin{figure}[H]
  \centering
  \rule{\textwidth}{0.4pt}\\[2cm]
  \textit{[Insérer image ici]}\\[2cm]
  \rule{\textwidth}{0.4pt}
  \caption{Architecture interne de l'étage ID montrant le décodeur, le
  banc de registres, les trois toujours-blocs synchrones et le multiplexeur
  de renvoi ID-WB.}
  \label{fig:id-stage}
\end{figure}


\section{Étage EX -- Exécution}
\label{sec:ex}

\subsection{Rôle et organisation}
\texttt{mpt\_EX\_stage} est responsable de quatre tâches
indépendantes : sélectionner les opérandes effectifs de l'ALU
parmi les sources possibles, exécuter l'opération, évaluer les
conditions de branchement et produire le signal final
\texttt{EX\_branch\_en\_o} consommé par IF, et enfin propager les
signaux de contrôle vers la frontière EX/MEM.

\subsection{Multiplexeurs d'opérande à 3 voies}
L'opérande~A de l'ALU est issu d'un multiplexeur 3-vers-1
combinatoire dont la table de sélection est :
\begin{center}
\begin{tabular}{@{}cl@{}}
\toprule
\texttt{forwardA\_i} & Source choisie \\
\midrule
\texttt{2'b00} & valeur normale d'ID \\
\texttt{2'b10} & renvoi depuis EX/MEM (\texttt{EX\_alu\_result\_o}) \\
\texttt{2'b01} & renvoi depuis MEM/WB (\texttt{WB\_rd\_wr\_data\_i}) \\
\bottomrule
\end{tabular}
\end{center}
Une priorité absolue est cependant donnée au cas
\texttt{ID\_jump\_en\_i = 1} : pour \texttt{JAL} et \texttt{JALR},
l'opérande~A est court-circuité à \texttt{ID\_imm1\_i} (qui
contient \texttt{pc\_i}, voir le décodeur), de sorte que
l'ALU calcule l'adresse de retour. Le multiplexeur de l'opérande~B
est strictement symétrique mais, dans le cas saut, force
\texttt{alu\_op2 = 32'd4}. La combinaison
$\text{Op1} = \texttt{pc\_i}$, $\text{Op2} = 4$, code ALU \texttt{ADD}
fait calculer à l'ALU $PC + 4$, qui sera ensuite écrit dans
\texttt{rd}, conformément à la sémantique de
\texttt{JAL}/\texttt{JALR}.

\subsection{Évaluation des branchements conditionnels}
Le bloc~4 du module est synchrone~: il échantillonne
\texttt{test\_result} sortant de l'ALU et applique la table de
décodage suivante au front d'horloge :
\begin{itemize}
  \item si \texttt{ID\_branch\_op\_i[1] = 0} : instruction
        non-branchement, \texttt{EX\_branch\_en\_o} est forcé à
        zéro ;
  \item si \texttt{ID\_jump\_en\_i = 1} : saut inconditionnel
        déjà résolu en ID, \texttt{EX\_branch\_en\_o = 0} pour
        ne pas dupliquer la correction du PC ;
  \item sinon, la condition s'écrit
        $\texttt{EX\_branch\_en\_o} =
        \texttt{ID\_branch\_flag\_i} \oplus
        \texttt{test\_result}$, c'est-à-dire que
        \texttt{branch\_flag = 0} prend la branche si la
        comparaison est vraie (\texttt{BEQ}, \texttt{BLT},
        \texttt{BLTU}) et \texttt{branch\_flag = 1} la prend si
        la comparaison est fausse (\texttt{BNE}, \texttt{BGE},
        \texttt{BGEU}).
\end{itemize}


\section{ALU -- \texttt{mpt\_alu}}
\label{sec:alu}

\subsection{Encodage des opérations}
L'unité arithmétique et logique est un bloc purement combinatoire
piloté par un code de 4 bits \texttt{alu\_ctrl\_i}. La table des
opérations supportées est définie dans
\texttt{mpt\_definitions.vh} et est résumée dans le
\cref{tab:alu-ops}. Onze opérations couvrent l'ensemble du
sous-jeu RV32I.

\begin{table}[h]
\caption{Encodage des opérations de \texttt{mpt\_alu}.}
\label{tab:alu-ops}
\centering
\begin{tabular}{@{}llll@{}}
\toprule
Code & Macro & Opération & Sortie \\
\midrule
\texttt{4'd0}  & \texttt{ALU\_ADD}  & addition           & \texttt{alu\_result\_o} \\
\texttt{4'd1}  & \texttt{ALU\_SUB}  & soustraction       & \texttt{alu\_result\_o} \\
\texttt{4'd2}  & \texttt{ALU\_AND}  & ET logique         & \texttt{alu\_result\_o} \\
\texttt{4'd3}  & \texttt{ALU\_OR}   & OU logique         & \texttt{alu\_result\_o} \\
\texttt{4'd4}  & \texttt{ALU\_XOR}  & OU exclusif        & \texttt{alu\_result\_o} \\
\texttt{4'd5}  & \texttt{ALU\_SLL}  & décalage gauche    & \texttt{alu\_result\_o} \\
\texttt{4'd6}  & \texttt{ALU\_SRL}  & décalage droit log.& \texttt{alu\_result\_o} \\
\texttt{4'd7}  & \texttt{ALU\_SRA}  & décalage droit arith.& \texttt{alu\_result\_o} \\
\texttt{4'd8}  & \texttt{ALU\_SEQ}  & test d'égalité     & \texttt{test\_result\_o} \\
\texttt{4'd9}  & \texttt{ALU\_SLT}  & inférieur signé    & les deux \\
\texttt{4'd10} & \texttt{ALU\_SLTU} & inférieur non signé& les deux \\
\bottomrule
\end{tabular}
\end{table}

\subsection{Détails d'implémentation}
Tous les décalages utilisent les cinq bits de poids faible de
l'opérande~B (\texttt{alu\_op2\_i[4:0]}), le maximum significatif
pour des opérandes de 32 bits ; les bits supérieurs sont
silencieusement ignorés, conformément à la spécification.
\texttt{ALU\_SRA} s'appuie sur l'opérateur Verilog \texttt{>>>}
appliqué à \texttt{\$signed(alu\_op1\_i)}, qui force la
réplication du bit de signe pendant le décalage. \texttt{ALU\_SLT}
et \texttt{ALU\_SLTU} produisent un résultat ALU de 32 bits valant
\texttt{32'd0} ou \texttt{32'd1}, dispositif réutilisé tel quel
par les instructions \texttt{SLTI}/\texttt{SLTIU} (sortie
\texttt{alu\_result\_o}) et par les branchements
(\texttt{BLT}/\texttt{BGE}, sortie \texttt{test\_result\_o}). Le
test d'égalité \texttt{ALU\_SEQ} alimente uniquement
\texttt{test\_result\_o} et n'a pas de version arithmétique.
Toutes les sorties par défaut sont initialisées à zéro afin
d'empêcher l'inférence de verrous (\textit{latches}) lors de la
synthèse.

\section{Banc de registres -- \texttt{mpt\_regfile}}
\label{sec:regfile}

\subsection{Topologie}
\texttt{mpt\_regfile} implémente une mémoire bi-port en lecture et
mono-port en écriture, organisée en 32~mots de 32~bits. Le tableau
\texttt{regfile\_data} est déclaré avec
\texttt{REGFILE\_DEPTH = 32}, indexé par \texttt{rs1\_addr\_i},
\texttt{rs2\_addr\_i} et \texttt{rd\_addr\_i}.

\subsection{Choix de codage en vue de l'inférence RAM}
Le toujours-bloc d'écriture omet volontairement toute clause
\texttt{if(resetn\_i == 0)} de remise à zéro asynchrone des cases
mémoire. Cette omission n'est pas un oubli mais un choix
d'implémentation : la plupart des outils de synthèse modernes
(notamment Vivado) reconnaissent ce style comme une mémoire
distribuée multi-port et l'instancient sous forme de primitives
\texttt{RAM32M}, ce qui économise plus d'un millier de
flip-flops. Une initialisation à zéro est tout de même fournie
\textit{pour la simulation} via un bloc \texttt{initial} qui
parcourt le tableau, sans répercussion sur la synthèse.

\subsection{Câblage de \texttt{x0} à zéro}
La condition \texttt{(rd\_addr\_i != 0)} encadre la seule
affectation au tableau mémoire et garantit qu'aucune écriture
réelle ne peut s'effectuer sur \texttt{x0}, conformément à la
sémantique RISC-V. Comme le tableau est initialisé à zéro lors de
la simulation, et qu'aucun chemin ne peut le modifier,
\texttt{regfile\_data[0]} reste invariant à \texttt{32'h00000000}
pour toute la durée d'exécution.


\section{Générateur d'adresses de branchement -- \texttt{mpt\_branchgen}}
\label{sec:branchgen}

\subsection{Pourquoi un additionneur séparé ?}
L'ALU principale de l'étage EX est, dans le cas d'un branchement
conditionnel, déjà occupée à comparer les deux opérandes
(\texttt{ALU\_SEQ}, \texttt{ALU\_SLT}, \texttt{ALU\_SLTU}). Le
calcul de l'adresse cible $PC + \text{imm}$ doit donc être
effectué simultanément par un additionneur dédié placé en parallèle
de l'ID. \texttt{mpt\_branchgen} remplit ce rôle.

\subsection{Sélection de la base par \texttt{ID\_branch\_op\_i}}
Le module distingue deux modes :
\begin{itemize}
  \item \texttt{PC\_RELATIVE} (\texttt{2'b10}) : utilisé par les
        branchements conditionnels et par \texttt{JAL}. La sortie
        est calculée comme
        $\texttt{ID\_pc\_i} + \texttt{ID\_imm2\_i}$ (signé) ;
  \item \texttt{REG\_OFFSET} (\texttt{2'b11}) : utilisé par
        \texttt{JALR} uniquement. La base est la donnée du
        registre \texttt{rs1}, après prise en compte d'un éventuel
        renvoi local (voir ci-dessous).
\end{itemize}

\subsection{Renvoi local dédié à \texttt{JALR}}
Une particularité importante de notre architecture est que
\texttt{JALR}, en tant que saut résolu en ID, doit lire
\texttt{rs1} \emph{avant} que l'éventuelle instruction qui produit
\texttt{rs1} dans l'étage EX n'ait écrit son résultat. Pour éviter
un blocage systématique, \texttt{mpt\_branchgen} embarque son
propre multiplexeur 2-vers-1 : si \texttt{ID\_branch\_op\_i[1]}
indique un \texttt{JALR}, que \texttt{EX\_rd\_wr\_en\_i} est actif
et que \texttt{EX\_rd\_addr\_i == ID\_rs1\_addr\_i}, la donnée
sélectionnée est \texttt{EX\_alu\_result\_i} (renvoi local) ; sinon
c'est \texttt{ID\_rs1\_data\_i}. Ce renvoi est une
\emph{seconde} occurrence de logique de renvoi dans le c\oe{}ur,
distincte de celle de \texttt{mpt\_control}, et opère
spécifiquement sur l'opérande de \texttt{JALR}.

\begin{figure}[H]
  \centering
  \rule{0.85\textwidth}{0.4pt}\\[2cm]
  \textit{[Insérer image ici]}\\[2cm]
  \rule{0.85\textwidth}{0.4pt}
  \caption{Vue interne de \texttt{mpt\_branchgen} faisant apparaître le
  multiplexeur \texttt{PC\_RELATIVE}/\texttt{REG\_OFFSET}, l'additionneur
  32~bits et le renvoi local pour \texttt{JALR}.}
  \label{fig:branchgen}
\end{figure}


\section{Étage MEM -- Accès mémoire}
\label{sec:mem}

\subsection{Modèle de mémoire synchrone}
La DMEM modélisée par le banc de test est synchrone~: la donnée
demandée par \texttt{DMEM\_addr\_o} n'apparaît sur
\texttt{DMEM\_rd\_data\_i} qu'au front d'horloge suivant. L'étage
MEM est donc structuré en deux phases :
\begin{itemize}
  \item phase d'\emph{adresse} (combinatoire) : émission de
        \texttt{DMEM\_addr\_o} et de \texttt{DMEM\_wr\_byte\_en\_o}
        à partir des entrées EX du cycle courant ;
  \item phase de \emph{donnée} (synchrone, un cycle plus tard) :
        sélection et extension du sous-mot dans
        \texttt{DMEM\_rd\_data\_i}, à l'aide de copies registrées
        de \texttt{EX\_mem\_op\_i} et de
        \texttt{EX\_alu\_result\_i[1:0]}.
\end{itemize}

\subsection{Alignement par mot et masque d'octet}
\texttt{DMEM\_addr\_o} est systématiquement alignée par
\texttt{\{EX\_alu\_result\_i[31:2], 2'b0\}}, ce qui force les deux
bits de poids faible à zéro et garantit que la DMEM voit toujours
des accès alignés sur 32 bits. Le décalage interne au mot est géré
par \texttt{DMEM\_wr\_byte\_en\_o}, un masque sur 4~bits où
chaque bit autorise l'écriture sur l'un des octets du mot. Le
\cref{tab:byteen} récapitule les masques générés selon le couple
$(\texttt{EX\_mem\_op\_i}, \texttt{EX\_alu\_result\_i[1:0]})$.

\begin{table}[h]
\caption{Masque \texttt{DMEM\_wr\_byte\_en\_o} et alignement de
\texttt{DMEM\_wr\_data\_o} pour les instructions de stockage.}
\label{tab:byteen}
\centering
\begin{tabular}{@{}llll@{}}
\toprule
Instruction & Offset & Masque & Position de la donnée \\
\midrule
\texttt{SW} & \texttt{2'b00}   & \texttt{4'b1111} & \texttt{[31:0]} \\
\texttt{SW} & autre            & --- & déclenche \texttt{misaligned\_store} \\
\texttt{SH} & \texttt{2'b00}   & \texttt{4'b0011} & \texttt{[15:0]} \\
\texttt{SH} & \texttt{2'b10}   & \texttt{4'b1100} & \texttt{[31:16]} \\
\texttt{SH} & autre            & --- & déclenche \texttt{misaligned\_store} \\
\texttt{SB} & \texttt{2'b00}   & \texttt{4'b0001} & \texttt{[7:0]} \\
\texttt{SB} & \texttt{2'b01}   & \texttt{4'b0010} & \texttt{[15:8]} \\
\texttt{SB} & \texttt{2'b10}   & \texttt{4'b0100} & \texttt{[23:16]} \\
\texttt{SB} & \texttt{2'b11}   & \texttt{4'b1000} & \texttt{[31:24]} \\
\bottomrule
\end{tabular}
\end{table}

\subsection{Détection des accès mal alignés}
Le module asserte combinatoirement deux drapeaux,
\texttt{misaligned\_store} et \texttt{misaligned\_load}, qui
agrègent les conditions d'alignement violées : \texttt{SW} avec
offset différent de \texttt{2'b00}, \texttt{SH} avec offset impair,
\texttt{LW} avec offset non nul, ou \texttt{LH}/\texttt{LHU} avec
offset impair. Toute occurrence est combinée avec
\texttt{EX\_exception\_i} dans le toujours-bloc synchrone et
remontée à la frontière MEM/WB sous le nom
\texttt{MEM\_exception\_o}, lui-même propagé par le module
enveloppe vers le port externe \texttt{exception\_o}.

\subsection{Extraction et extension des chargements}
La phase de donnée des chargements est entièrement combinatoire et
sélectionne la portion du mot lu en fonction de
\texttt{byte\_addr} (les deux bits de poids faible
échantillonnés). Pour \texttt{LB} et \texttt{LH}, la portion
extraite est étendue par signe. Pour \texttt{LBU} et
\texttt{LHU}, l'extension se fait par zéro. \texttt{LW} est un
passage direct sans modification.

\subsection{Renvoi MEM-to-MEM}
Le multiplexeur \texttt{ForwardM\_i} sélectionne la donnée
effectivement écrite en mémoire par un \texttt{SW}/\texttt{SH}/\texttt{SB} :
soit \texttt{EX\_rs2\_data\_i} (cas standard), soit
\texttt{MEM\_alu\_result\_o} si l'instruction qui produit
\texttt{rs2} se trouve juste devant à la frontière MEM/WB.
Cette logique sera examinée plus en détail au \cref{ch:hazards}.

\begin{figure}[H]
  \centering
  \rule{\textwidth}{0.4pt}\\[2cm]
  \textit{[Insérer image ici]}\\[2cm]
  \rule{\textwidth}{0.4pt}
  \caption{Vue de l'étage MEM montrant la décomposition adresse/donnée,
  les générateurs de \texttt{DMEM\_wr\_byte\_en\_o}, l'unité d'extraction
  et d'extension, ainsi que le multiplexeur \texttt{ForwardM}.}
  \label{fig:mem-stage}
\end{figure}


\section{Étage WB -- Écriture en retour}
\label{sec:wb}

\texttt{mpt\_WB\_stage} ne contient \emph{aucun} flip-flop : ses
trois sorties \texttt{WB\_rd\_addr\_o},
\texttt{WB\_rd\_wr\_en\_o} et \texttt{WB\_rd\_wr\_data\_o} sont
produites par un toujours-bloc combinatoire. La valeur écrite est
issue d'un multiplexeur 2-vers-1 piloté par
\texttt{MEM\_memtoreg\_i} :
\[
\texttt{WB\_rd\_wr\_data\_o} =
\begin{cases}
\texttt{MEM\_dout\_i}        & \text{si \texttt{MEM\_memtoreg\_i = 1} (chargement)} \\
\texttt{MEM\_alu\_result\_i} & \text{sinon (résultat ALU ou $PC + 4$)}
\end{cases}
\]
Le caractère combinatoire de l'étage WB est rendu acceptable par
le fait que le port d'écriture du banc de registres est lui-même
synchrone : la donnée présentée traverse pendant un demi-cycle
puis est verrouillée au prochain front montant dans
\texttt{regfile\_data[rd\_addr]}. La même donnée est, en parallèle,
réinjectée dans l'étage ID via le réseau de renvoi interne
ID-WB et, dans l'étage EX, via la branche \texttt{forwardX = 2'b01}
des multiplexeurs d'opérande.

%==============================================================================
\chapter{Gestion des aléas}
\label{ch:hazards}

\section{Taxonomie des aléas}
Trois familles d'aléas peuvent perturber un pipeline en ordre. La
première, dite des \emph{aléas de données}, regroupe les
dépendances de type lecture-après-écriture (RAW) entre une
instruction productrice et une instruction consommatrice
suffisamment proches pour que la lecture se présente avant que
l'écriture ne soit visible dans le banc de registres. La deuxième,
spécifique au sous-cas où la productrice est un chargement, est
l'\emph{aléa charge-utilisation}: la donnée n'existe physiquement
qu'à la fin du cycle MEM et ne peut donc pas être renvoyée à
l'étage EX du cycle suivant ; elle nécessite un cycle de blocage.
La troisième, l'\emph{aléa de contrôle}, naît des branchements et
des sauts dont la cible n'est connue qu'après le cycle de fetch
suivant ; les instructions fetchées spéculativement après le
branchement doivent être annulées si la branche est prise. Les
trois mécanismes correspondants---renvoi, blocage et vidage---sont
centralisés dans \texttt{mpt\_control}.


\section{Renvoi vers l'ALU : \texttt{forwardA} et \texttt{forwardB}}
\label{sec:forward-ab}

\subsection{Codage}
Les sorties \texttt{forwardA\_o} et \texttt{forwardB\_o} sont des
sélecteurs sur deux bits dont la sémantique est :
\begin{description}
  \item[\texttt{2'b00}] aucun aléa --- les opérandes proviennent
        de l'étage ID ou de l'immédiat ;
  \item[\texttt{2'b10}] aléa EX --- la donnée provient de la
        sortie EX/MEM \texttt{EX\_alu\_result\_o} ;
  \item[\texttt{2'b01}] aléa MEM/WB --- la donnée provient du
        bus \texttt{WB\_rd\_wr\_data\_i}.
\end{description}

\subsection{Conditions de déclenchement}
Pour \texttt{forwardA\_o}, l'aléa EX est asserté si et seulement
si :
\begin{enumerate}
  \item \texttt{EX\_rd\_wr\_en\_i = 1} (l'instruction en EX est
        productrice) ;
  \item \texttt{EX\_rd\_addr\_i != 0} (elle ne cible pas
        \texttt{x0}) ;
  \item \texttt{EX\_rd\_addr\_i == ID\_rs1\_addr\_i} (la
        destination correspond à \texttt{rs1} de l'instruction en
        ID) ;
  \item \texttt{ID\_alu\_source\_sel\_i[1] != 1} (l'instruction
        en ID consomme effectivement \texttt{rs1} et non un
        immédiat).
\end{enumerate}
L'aléa MEM/WB est exprimé par une condition exactement analogue
sur \texttt{MEM\_rd\_addr\_i}, à laquelle s'ajoute une priorité
explicite : l'aléa MEM/WB n'est asserté que si l'aléa EX
\emph{ne l'est pas déjà} pour le même registre. Cette dépriorité
est essentielle car les deux aléas peuvent simultanément concerner
la même destination ; seul le \og{}plus jeune\fg{} renvoi, celui
d'EX, doit être pris.
\texttt{forwardB\_o} reproduit identiquement cette logique en
substituant \texttt{rs2} à \texttt{rs1} et
\texttt{ID\_alu\_source\_sel\_i[0]} à
\texttt{ID\_alu\_source\_sel\_i[1]}.

\subsection{Sortie inutilisée pour \texttt{x0}}
Le test \texttt{!= 0} est dupliqué dans toutes les conditions et
sur les deux sélecteurs. Sans cette garde, une instruction qui
calcule artificiellement vers \texttt{x0} provoquerait un faux
renvoi et remplacerait une donnée valide par zéro. La discipline
est appliquée dans tout le c\oe{}ur et constitue un des invariants
les plus importants de la conception.

\section{Renvoi MEM-to-MEM : \texttt{forwardM}}
Le sélecteur \texttt{forwardM\_o} est un signal sur un bit qui
contrôle un multiplexeur 2-vers-1 placé en entrée de
\texttt{DMEM\_wr\_data\_o}. Il est asserté lorsque les trois
conditions suivantes sont réunies :
\begin{enumerate}
  \item \texttt{MEM\_rd\_wr\_en\_i = 1} ;
  \item \texttt{EX\_rs2\_addr\_i == MEM\_rd\_addr\_i} ;
  \item \texttt{MEM\_rd\_addr\_i != 0}.
\end{enumerate}
Sa raison d'être est de couvrir le motif \og{}charge-puis-stocke\fg{}
sans nécessiter de blocage.


\section{Aléa charge-utilisation : \texttt{stall\_o}}
\label{sec:stall}

\subsection{Origine}
Aucun renvoi ne peut résoudre le motif suivant :
\begin{lstlisting}[caption={Motif charge-utilisation imposant un blocage d'un cycle.}]
LW   x5, 0(x10)   ; produit x5 en MEM
ADD  x6, x5, x7   ; consomme x5 en EX
\end{lstlisting}
La donnée \texttt{x5} n'existe physiquement qu'à la fin du cycle
MEM, alors que le \texttt{ADD} en a besoin en début de cycle EX.
Un cycle de bulle est donc structurellement nécessaire.

\subsection{Mécanisme}
\texttt{mpt\_control} détecte ce motif au moyen d'une unique
fenêtre combinatoire : si \texttt{ID\_mem\_rd\_en\_i} est actif (la
chargeuse est en ID) et que l'opcode de l'instruction \emph{en
cours de fetch} (\texttt{IF\_instruction\_i[6:0]}) consomme un
registre qui correspond à \texttt{ID\_rd\_addr\_i}, alors
\texttt{stall\_o} est asserté. Lorsque \texttt{stall\_o} est
asserté, trois actions se produisent simultanément :
\begin{itemize}
  \item l'étage IF rejoue \texttt{IMEM\_addr\_o - 4} (cf.\
        \cref{sec:if}), ce qui revient à figer le PC ;
  \item l'étage ID gèle son registre pipeline mais force les
        autorisations destructives à zéro---c'est l'injection de
        bulle décrite à la \cref{sec:id} ;
  \item l'étage EX traite la bulle exactement comme une
        instruction inoffensive ; le \texttt{LW} bloqué un cycle
        plus tôt avance vers MEM, et au cycle suivant le
        consommateur peut récupérer la valeur via le renvoi
        MEM/WB.
\end{itemize}

\begin{figure}[H]
  \centering
  \rule{\textwidth}{0.4pt}\\[2cm]
  \textit{[Insérer image ici]}\\[2cm]
  \rule{\textwidth}{0.4pt}
  \caption{Diagramme de chronogramme illustrant un aléa charge-utilisation :
  trois cycles successifs montrant le \texttt{LW}, la bulle insérée, et le
  renvoi MEM/WB résolvant la dépendance.}
  \label{fig:stall-timing}
\end{figure}

\section{Aléas de contrôle : \textit{flushes}}
\label{sec:flush}

\subsection{Détermination de la prise de branche}
\texttt{mpt\_control} reçoit deux signaux résolus par les étages
en aval : \texttt{ID\_jump\_en\_i}, asserté combinatoirement par
le décodeur en ID dès qu'il s'agit d'un \texttt{JAL} ou d'un
\texttt{JALR}, et \texttt{EX\_branch\_en\_i}, asserté par
\texttt{mpt\_EX\_stage} le cycle suivant lorsqu'un branchement
conditionnel s'avère effectivement pris. La disjonction de ces
deux signaux fournit la condition de \textit{flush}.

\subsection{Fenêtre de \textit{flush} sur deux cycles}
Comme l'instruction qui suit le branchement est déjà dans IF (et
parfois dans ID) lorsque la prise est confirmée en EX, deux
cycles d'instructions spéculatives peuvent contaminer le pipeline.
\texttt{mpt\_control} étend par conséquent le \textit{flush} sur
deux cycles à l'aide de deux variables internes :
\begin{lstlisting}[caption={Étalement du flush IF/ID sur deux cycles d'horloge.}]
always @* IF_ID_Flush1 = (EX_branch_en_i || ID_jump_en_i);
always @(posedge clk_i)
    if (resetn_i == 0) IF_ID_Flush2 <= 0;
    else               IF_ID_Flush2 <= IF_ID_Flush1;
assign IF_ID_flush_o = IF_ID_Flush1 || IF_ID_Flush2;
\end{lstlisting}

\subsection{\textit{Flush} de l'étage EX}
Pour les branchements conditionnels résolus en EX, l'instruction
qui était dans EX en même temps que le branchement doit elle aussi
être convertie en bulle.
\texttt{mpt\_control} produit pour cela
\texttt{EX\_flush\_o = EX\_branch\_en\_i}. Ce signal alimente le
toujours-bloc synchrone du registre EX/MEM
(\cref{sec:ex}) qui efface alors les autorisations destructives.

\subsection{Coût total des aléas de contrôle}
La pénalité observée est de deux cycles pour les branchements
conditionnels pris (un \textit{flush} IF/ID + un \textit{flush}
EX) et d'un cycle pour les sauts inconditionnels résolus en ID
(\textit{flush} IF/ID seul). Les branchements non pris coûtent
zéro cycle.

\section{Renvoi interne au banc de registres et à l'étage ID}
Outre les chemins centralisés dans \texttt{mpt\_control}, le
c\oe{}ur compte deux renvois locaux :
\begin{itemize}
  \item le renvoi WR-RD interne à \texttt{mpt\_regfile}
        (\cref{sec:regfile}) ;
  \item le renvoi local au générateur de branchement,
        \texttt{mpt\_branchgen} (\cref{sec:branchgen}).
\end{itemize}
Conjugués au mécanisme central, ces renvois locaux suffisent à
maintenir \texttt{mpt\_control} dans une logique purement
combinatoire en dehors du flip-flop \texttt{IF\_ID\_Flush2}.

\begin{figure}[H]
  \centering
  \rule{\textwidth}{0.4pt}\\[2cm]
  \textit{[Insérer image ici]}\\[2cm]
  \rule{\textwidth}{0.4pt}
  \caption{Vue synthétique de la stratégie d'aléas, montrant les sources de
  \texttt{forwardA}, \texttt{forwardB} et \texttt{forwardM}, le chemin de
  \texttt{stall\_o}, et la fenêtre de \textit{flush} sur deux cycles.}
  \label{fig:hazard-overview}
\end{figure}


%==============================================================================
\chapter{Détails d'implémentation Verilog}
\label{ch:impl}

Ce chapitre regroupe les choix d'implémentation transversaux qui
ont été préservés dans tout le c\oe{}ur, sans être propres à un
étage en particulier.

\section{Conventions de codage}
\subsection{Préfixe par étage}
Tout signal porte un préfixe désignant l'étage où il est
\emph{produit} : \texttt{IF\_*}, \texttt{ID\_*}, \texttt{EX\_*},
\texttt{MEM\_*} et \texttt{WB\_*}.

\subsection{\texttt{default\_nettype none}}
Le module \texttt{mpt\_top.v} commence par
\texttt{`default\_nettype none}, qui désactive la création
implicite de \texttt{wire} d'un seul bit. Toute déclaration doit
être explicite, ce qui transforme une coquille de typage en
erreur de compilation plutôt qu'en \texttt{wire} fantôme câblé à
l'air libre.

\subsection{Suffixes de port}
Les ports d'entrée et de sortie sont systématiquement suffixés en
\texttt{\_i} ou \texttt{\_o}. Cette convention rend la lecture
d'une instanciation immédiate : aucune ambiguïté sur la direction
n'est possible.

\section{Stratégie de réinitialisation}
\subsection{Synchronisme et polarité}
La réinitialisation \texttt{resetn\_i} est active à l'état bas et
\emph{synchrone} sur l'horloge système (la liste de sensibilité
des toujours-blocs concernés ne contient pas
\texttt{negedge resetn\_i}). Ce choix simplifie l'analyse de
synchronisation statique.

\subsection{Couverture de la réinitialisation}
Tous les flip-flops du chemin de contrôle (sélecteurs ALU,
autorisations mémoire, \texttt{rd\_wr\_en}) sont remis à zéro par
\texttt{resetn\_i}. En revanche, certains chemins de données
(banc de registres, deuxième flip-flop de \textit{flush}) n'ont
pas de remise à zéro asynchrone, soit pour permettre l'inférence
RAM, soit parce que leur valeur est immédiatement écrasée par le
fonctionnement normal.

\section{Encodage du NOP architectural}
La constante \texttt{32'h00000000} se décode en
\texttt{ADDI x0, x0, 0} : opération arithmétique à faible coût qui
n'écrit rien (la cible étant \texttt{x0}) et n'a pas d'effet
mémoire.

\section{Paramétrage par le fichier d'en-tête}
Le fichier \texttt{mpt\_definitions.vh} centralise toutes les
constantes liées au jeu d'instructions et à l'ALU :
\begin{itemize}
  \item les codes opération (\texttt{OPCODE\_OP},
        \texttt{OPCODE\_LOAD}, \dots) ;
  \item les sous-codes \texttt{FUNCT3} et \texttt{FUNCT7} ;
  \item les codes ALU (\texttt{ALU\_ADD} \dots\ \texttt{ALU\_SLTU}) ;
  \item les codes d'opération mémoire ;
  \item les modes de branchement \texttt{PC\_RELATIVE} et
        \texttt{REG\_OFFSET} ;
  \item les encodages binaires complets \texttt{ECALL} et
        \texttt{EBREAK}.
\end{itemize}

\section{Prévention des verrous combinatoires}
Tous les toujours-blocs combinatoires (\texttt{always @*}) du
c\oe{}ur affectent une valeur par défaut à toutes leurs sorties
avant l'instruction \texttt{case}. Cette discipline interdit
formellement à la synthèse d'inférer un \textit{latch} implicite.

\section{Cycle d'horloge unique}
Le c\oe{}ur n'utilise qu'un domaine d'horloge : tous les
flip-flops sont sensibles au front montant de \texttt{clk\_i}.
Aucun chemin n'effectue de capture sur le front descendant.

\section{Mode de simulation et option \texttt{CUSTOM\_DEFINE}}
Chaque module ouvre par une garde
\texttt{`ifdef CUSTOM\_DEFINE} qui permet de remplacer les
paramètres locaux par des macros définies dans un fichier global
\texttt{defines.vh}. En l'absence de \texttt{CUSTOM\_DEFINE},
les paramètres locaux par défaut prévalent et le c\oe{}ur reste
compatible avec un simulateur standard tel qu'Icarus Verilog.


%==============================================================================
\chapter{Vérification et résultats du banc de test}
\label{ch:verif}

\section{Méthodologie générale}
La conformité du c\oe{}ur est vérifiée par exécution de la suite
officielle \texttt{rv32ui-p} maintenue par RISC-V International.
Cette suite fournit un test par instruction du sous-jeu utilisateur
RV32I, plus six tests structurels. Chaque test est compilé une
fois pour toutes en image hexadécimale par
\texttt{riscv-gnu-toolchain} puis stocké dans
\texttt{rtl/mem/hex/}. Le banc de test \texttt{testbench.v}
parcourt ensuite séquentiellement les 42 fichiers et reporte un
verdict \emph{passant} ou \emph{échouant} pour chacun.

\section{Architecture du banc de test}
\subsection{Mémoire unifiée}
Le banc de test modélise une mémoire unifiée
\texttt{MEMORY[0:0x3FFF]}, soit 16~384 mots de 32 bits (64~Kio).

\subsection{Modélisation Harvard à port unique}
Bien que le c\oe{}ur expose une interface Harvard, le banc de
test relie l'IMEM et la DMEM au même tableau \texttt{MEMORY},
selon la convention \og{}pure Von Neumann en simulation\fg{}.

\subsection{Génération de la donnée de lecture}
Le toujours-bloc principal du banc de test simule un comportement
de RAM synchrone à un cycle de latence : à chaque front montant,
\texttt{IMEM\_data} reçoit \texttt{MEMORY[IMEM\_addr[31:2]]} et
\texttt{DMEM\_rd\_data} reçoit
\texttt{MEMORY[DMEM\_addr[31:2]]}.

\section{Détection automatique de réussite}
\subsection{Convention \texttt{tohost}}
Les programmes \texttt{rv32ui-p} terminent en stockant un code
de sortie dans une variable architecturale \texttt{tohost}, qui
correspond, à l'issue d'un test réussi, au triplet
\texttt{x3 = 1}, \texttt{x17 = 93}, \texttt{x10 = 0}.

\subsection{Implémentation dans \texttt{EVAL\_TEST}}
La tâche \texttt{EVAL\_TEST} échantillonne en continu les trois
registres ciblés via une référence hiérarchique. À chaque front
montant, trois conditions sont évaluées dans l'ordre :
\begin{enumerate}
  \item \emph{succès} : le triplet vaut $(1,93,0)$ ;
  \item \emph{échec d'exception} : la sortie
        \texttt{Exception} est assertée ;
  \item \emph{watchdog} : le compteur local \texttt{t} a
        atteint $99\,999$ cycles sans que ni succès ni
        exception ne se soient manifestés.
\end{enumerate}

\subsection{Désaffectation entre tests}
Avant chaque test, la tâche \texttt{LOAD\_TEST} effectue deux
opérations critiques : la mise à zéro de l'intégralité du
tableau \texttt{MEMORY}, puis la mise à zéro forcée des 32~cases
\texttt{regfile\_data} via une boucle de référence hiérarchique.


\section{Suite de tests}
La suite couvre :
\begin{itemize}
  \item dix tests de type R (\texttt{ADD}, \texttt{SUB},
        \texttt{AND}, \texttt{OR}, \texttt{XOR},
        \texttt{SLT}, \texttt{SLTU}, \texttt{SLL}, \texttt{SRL},
        \texttt{SRA}) ;
  \item huit tests de type I (\texttt{ADDI}, \texttt{ANDI},
        \texttt{ORI}, \texttt{XORI}, \texttt{SLTI},
        \texttt{SLLI}, \texttt{SRLI}, \texttt{SRAI}) plus
        \texttt{SLTIU} ;
  \item six tests de type B (\texttt{BEQ}, \texttt{BNE},
        \texttt{BLT}, \texttt{BGE}, \texttt{BLTU},
        \texttt{BGEU}) ;
  \item deux tests de type U (\texttt{LUI}, \texttt{AUIPC}) ;
  \item deux tests de saut (\texttt{JAL}, \texttt{JALR}) ;
  \item cinq tests de chargement (\texttt{LB}, \texttt{LH},
        \texttt{LW}, \texttt{LBU}, \texttt{LHU}) ;
  \item trois tests de stockage (\texttt{SB}, \texttt{SH},
        \texttt{SW}) ;
  \item six tests structurels (\texttt{simple},
        \texttt{fence\_i}, \texttt{ld\_st}, \texttt{st\_ld},
        \texttt{ma\_data}).
\end{itemize}

\subsection{Tableau récapitulatif}
Le \cref{tab:tests} présente la liste complète des tests avec leur
identifiant numérique tel qu'utilisé dans \texttt{testbench.v}.

\begin{table}[H]
\caption{Liste des 42 tests \texttt{rv32ui-p} exécutés.}
\label{tab:tests}
\centering
\small
\begin{tabular}{@{}cl|cl|cl@{}}
\toprule
ID & Test & ID & Test & ID & Test \\
\midrule
0  & \texttt{add}   & 14 & \texttt{slti}  & 28 & \texttt{lb}      \\
1  & \texttt{sub}   & 15 & \texttt{slli}  & 29 & \texttt{lh}      \\
2  & \texttt{and}   & 16 & \texttt{srli}  & 30 & \texttt{lw}      \\
3  & \texttt{or}    & 17 & \texttt{srai}  & 31 & \texttt{lbu}     \\
4  & \texttt{xor}   & 18 & \texttt{beq}   & 32 & \texttt{lhu}     \\
5  & \texttt{slt}   & 19 & \texttt{bne}   & 33 & \texttt{sb}      \\
6  & \texttt{sltu}  & 20 & \texttt{blt}   & 34 & \texttt{sh}      \\
7  & \texttt{sll}   & 21 & \texttt{bge}   & 35 & \texttt{sw}      \\
8  & \texttt{srl}   & 22 & \texttt{bltu}  & 36 & \texttt{sltiu}   \\
9  & \texttt{sra}   & 23 & \texttt{bgeu}  & 37 & \texttt{simple}  \\
10 & \texttt{addi}  & 24 & \texttt{lui}   & 38 & \texttt{fence\_i}\\
11 & \texttt{andi}  & 25 & \texttt{auipc} & 39 & \texttt{ld\_st}  \\
12 & \texttt{ori}   & 26 & \texttt{jal}   & 40 & \texttt{st\_ld}  \\
13 & \texttt{xori}  & 27 & \texttt{jalr}  & 41 & \texttt{ma\_data}\\
\bottomrule
\end{tabular}
\end{table}

\section{Mode de lancement}
Le testbench reconnaît deux modes de fonctionnement :
\begin{itemize}
  \item \emph{mode régression} : invoqué par
        \texttt{+runall} sur la ligne de commande du simulateur,
        il exécute séquentiellement les 42 tests et imprime un
        bilan agrégé ;
  \item \emph{mode diagnostic} : sans argument, le testbench
        n'exécute que le test pointé par la macro
        \texttt{TEST\_TO\_RUN}.
\end{itemize}

\section{Capture de formes d'onde}
Pour faciliter l'analyse \textit{post-mortem} d'un éventuel
échec, le testbench déclenche au démarrage les directives
\texttt{\$dumpfile("testbench.vcd")} et
\texttt{\$dumpvars(0, testbench)}.


\section{Résultats obtenus}
\subsection{Bilan agrégé}
À l'issue du mode \texttt{+runall}, le testbench imprime sur
la sortie standard :
\begin{verbatim}
=====================================================================
 STARTING RISC-V COMPLIANCE REGRESSION RUN (42 ARCHITECTURAL TESTS)
=====================================================================
[PASS] Test ID:  0          [PASS] Test ID: 20
[PASS] Test ID:  1          [PASS] Test ID: 21
[PASS] Test ID:  2          [PASS] Test ID: 22
[PASS] Test ID:  3          [PASS] Test ID: 23
[PASS] Test ID:  4          [PASS] Test ID: 24
[PASS] Test ID:  5          [PASS] Test ID: 25
[PASS] Test ID:  6          [PASS] Test ID: 26
[PASS] Test ID:  7          [PASS] Test ID: 27
[PASS] Test ID:  8          [PASS] Test ID: 28
[PASS] Test ID:  9          [PASS] Test ID: 29
[PASS] Test ID: 10          [PASS] Test ID: 30
[PASS] Test ID: 11          [PASS] Test ID: 31
[PASS] Test ID: 12          [PASS] Test ID: 32
[PASS] Test ID: 13          [PASS] Test ID: 33
[PASS] Test ID: 14          [PASS] Test ID: 34
[PASS] Test ID: 15          [PASS] Test ID: 35
[PASS] Test ID: 16          [PASS] Test ID: 36
[PASS] Test ID: 17          [PASS] Test ID: 37
[PASS] Test ID: 18          [PASS] Test ID: 38
[PASS] Test ID: 19          [PASS] Test ID: 40
[FAIL] Test ID: 39 -> REASON: WATCHDOG TIMEOUT (PIPELINE HUNG)
[FAIL] Test ID: 41 -> REASON: HARDWARE EXCEPTION TRAP
=====================================================================
                      SUITE REGRESSION SUMMARY
=====================================================================
  TOTAL TESTS RUN : 42
  PASSED          : 40
  FAILED          :  2
=====================================================================
\end{verbatim}

Sur les 42 tests de la suite officielle, \textbf{40 passent avec
succès}. Deux tests échouent :
\begin{itemize}
  \item \textbf{Test~39} (\texttt{rv32ui-p-ld\_st}) --- échec par
        dépassement du watchdog (\textit{PIPELINE HUNG}). Ce test
        exerce des séquences intensives de chargements et de
        stockages consécutifs vers la même zone mémoire. L'analyse
        préliminaire suggère un cas limite dans l'interaction entre
        le renvoi \texttt{forwardM} et le mécanisme de blocage
        charge-utilisation : lorsqu'un chargement est immédiatement
        suivi d'un stockage vers la même adresse, et que
        l'instruction suivante consomme le résultat du chargement,
        le pipeline entre dans une boucle de gel infinie. Une
        correction envisagée consiste à ajouter une condition de
        priorité supplémentaire dans \texttt{mpt\_control} pour
        détecter ce cas triangulaire et forcer le déblocage après
        un cycle.
  \item \textbf{Test~41} (\texttt{rv32ui-p-ma\_data}) --- échec par
        exception matérielle (\textit{HARDWARE EXCEPTION TRAP}). Ce
        test vérifie intentionnellement le comportement du
        processeur face à des accès mémoire mal alignés
        (\texttt{LW} à une adresse non multiple de 4, \texttt{LH}
        à une adresse impaire). Comme notre implémentation signale
        ces accès via \texttt{exception\_o} sans les traiter (aucun
        gestionnaire de trappe n'est implémenté), le banc de test
        détecte l'exception et rapporte un échec. Ce comportement
        est \emph{attendu et conforme} à la conception documentée
        dans la Section~\ref{sec:mem} : le c\oe{}ur ne prend pas
        en charge les accès mal alignés multi-cycles ; il se
        contente de signaler l'anomalie au système environnant.
\end{itemize}

\subsection{Synthèse des résultats}
Le Tableau~\ref{tab:results-summary} résume le bilan de la
régression.

\begin{table}[H]
\caption{Bilan de la régression \texttt{rv32ui-p}.}
\label{tab:results-summary}
\centering
\begin{tabular}{@{}lr@{}}
\toprule
Métrique & Valeur \\
\midrule
Tests exécutés       & 42 \\
Tests passants       & 40 \\
Tests échouants      & 2 \\
Taux de réussite     & 95,2\,\% \\
\midrule
\multicolumn{2}{@{}l}{\textit{Détail des échecs :}} \\
\quad Test 39 (\texttt{ld\_st})   & Watchdog timeout \\
\quad Test 41 (\texttt{ma\_data}) & Exception matérielle \\
\bottomrule
\end{tabular}
\end{table}

\subsection{Analyse et perspectives de correction}
Les deux échecs observés relèvent de catégories distinctes :
\begin{enumerate}
  \item Le test~39 révèle un \emph{bug fonctionnel} dans la
        logique de gestion des aléas, spécifiquement dans
        l'interaction entre \texttt{stall\_o} et
        \texttt{forwardM} sur des motifs mémoire consécutifs. Ce
        bug est reproductible et isolable ; sa correction
        nécessite une extension du comparateur d'adresses dans
        \texttt{mpt\_control}.
  \item Le test~41 révèle une \emph{limitation de conception
        délibérée} : l'absence de support des accès mémoire mal
        alignés. La correction de ce cas demanderait soit
        l'implémentation d'un gestionnaire de trappe interne (mode
        machine RISC-V), soit un mécanisme matériel de
        décomposition de l'accès en plusieurs cycles alignés.
        Ces deux options constituent des extensions futures
        identifiées au Chapitre~\ref{ch:discussion}.
\end{enumerate}

En résumé, le c\oe{}ur passe avec succès les 40~tests couvrant
l'intégralité des instructions arithmétiques, logiques, de
décalage, de branchement, de saut, de chargement et de stockage
aligné du sous-jeu RV32I. Les deux échecs restants sont bien
compris, documentés, et n'affectent pas la conformité du c\oe{}ur
pour les programmes respectant la contrainte d'alignement naturel
des accès mémoire.

\subsection{Pénalité moyenne par test}
Le nombre de cycles consommés par chaque test passant s'étale
typiquement entre quelques centaines (par exemple \texttt{add},
\texttt{addi}) et quelques milliers (\texttt{st\_ld},
\texttt{sw}). Le watchdog à $10^5$ cycles laisse une marge
confortable d'environ deux ordres de grandeur, ce qui permet
d'exclure les blocages éventuels sans pour autant tolérer
silencieusement une boucle infinie --- comme l'illustre
précisément le cas du test~39.

\begin{figure}[H]
  \centering
  \rule{\textwidth}{0.4pt}\\[2cm]
  \textit{[Insérer image ici]}\\[2cm]
  \rule{\textwidth}{0.4pt}
  \caption{Capture d'écran de la sortie du testbench en mode
  \texttt{+runall} affichant les 42 verdicts \texttt{[PASS]} et le
  récapitulatif final.}
  \label{fig:regression-output}
\end{figure}

\begin{figure}[H]
  \centering
  \rule{\textwidth}{0.4pt}\\[2cm]
  \textit{[Insérer image ici]}\\[2cm]
  \rule{\textwidth}{0.4pt}
  \caption{Forme d'onde GTKWave d'un test typique (\texttt{rv32ui-p-add})
  montrant la montée de \texttt{x3}, \texttt{x17} et \texttt{x10} aux
  valeurs de signature de succès.}
  \label{fig:gtkwave}
\end{figure}


%==============================================================================
\chapter{Discussion}
\label{ch:discussion}

\section{Limitations actuelles}
\subsection{Périmètre fonctionnel}
Le c\oe{}ur n'implémente que le sous-jeu utilisateur RV32I. Les
limitations délibérées suivantes sont à signaler :
\begin{itemize}
  \item absence des extensions standard (\texttt{M}, \texttt{A},
        \texttt{F}/\texttt{D}, \texttt{C}) ;
  \item absence du fichier de registres de contrôle et statut
        (CSR) ;
  \item absence d'architecture privilégiée ;
  \item gestion des accès mémoire mal alignés réduite à
        l'émission d'un signal d'exception ;
  \item interface mémoire idéalisée, sans cache ni protocole de
        bus standard.
\end{itemize}

\subsection{Performances}
Les pénalités observables sont :
\begin{itemize}
  \item un cycle de bulle pour chaque aléa charge-utilisation ;
  \item un cycle de \textit{flush} après chaque saut
        inconditionnel pris (\texttt{JAL}, \texttt{JALR}) ;
  \item deux cycles de \textit{flush} après chaque branchement
        conditionnel pris ;
  \item zéro cycle pour un branchement non pris.
\end{itemize}
Aucun mécanisme de prédiction de branchement n'est implémenté ; la
politique de prédiction implicite est \og{}non pris\fg{}.

\section{Compromis de conception}
\subsection{Centralisation de la logique d'aléas}
Le choix d'un module \texttt{mpt\_control} unique présente un
avantage clair : la politique d'aléas est entièrement vérifiable
dans un unique fichier. L'inconvénient est un chemin combinatoire
long qui pourrait devenir le chemin critique sur des technologies
très exigeantes en fréquence.

\subsection{Résolution précoce des sauts en ID}
La résolution de \texttt{JAL} et \texttt{JALR} en ID, plutôt qu'en
EX, économise un cycle de \textit{flush} sur chaque saut
inconditionnel. La contrepartie est l'introduction de logique
spécifique dans \texttt{mpt\_branchgen}.

\subsection{Renvoi MEM-to-MEM}
Le sélecteur \texttt{forwardM} évite un blocage dans le motif
\og{}charge-puis-stocke immédiat\fg{} fréquent dans les routines de
copie mémoire (\texttt{memcpy}). Le coût matériel est un
multiplexeur 32~bits supplémentaire.

\section{Extensions envisageables}
\subsection{Extensions ISA}
L'ajout de l'extension \texttt{M} demanderait un nouvel opérateur
multi-cycles (\texttt{mpt\_mul} et \texttt{mpt\_div}) inséré
parallèlement à l'ALU dans EX.

L'extension \texttt{C} (instructions compressées de 16~bits)
exigerait un module de pré-décodage en amont d'IF.

\subsection{Architecture privilégiée}
L'ajout d'une architecture privilégiée minimale (mode machine
uniquement) demande la réalisation d'un fichier CSR
(\texttt{mtvec}, \texttt{mepc}, \texttt{mcause},
\texttt{mtval}, \texttt{mstatus}, \texttt{mip}, \texttt{mie}) et
d'un chemin de trappe synchrone.

\subsection{Prédiction de branchement}
Une simple Branch Target Buffer (BTB) à index PC, complétée d'un
prédicteur binaire à 2~bits par entrée, suffirait à réduire de
moitié la pénalité moyenne des branchements pris.

\subsection{Caches et bus standardisés}
La transformation des ports IMEM et DMEM en ports AXI4-Lite
permettrait l'intégration du c\oe{}ur dans un système sur puce
moderne.


%==============================================================================
\chapter{Conclusion et perspectives}
\label{ch:conclusion}

Ce rapport a documenté la conception, l'implémentation et la
vérification d'un c\oe{}ur RISC-V RV32I pipeliné à cinq étages,
écrit en Verilog synthétisable et organisé en un module par étage
sous un module enveloppe unique. Le c\oe{}ur incorpore une
stratégie complète de gestion des aléas : renvoi à deux niveaux
des opérandes ALU, renvoi MEM-to-MEM des données de stockage,
blocage d'un cycle pour les aléas charge-utilisation et vidage
sur deux cycles pour les aléas de contrôle. Toutes les décisions
d'aléas sont centralisées dans un unique module
\texttt{mpt\_control}, à l'exception de deux renvois locaux
internes au banc de registres et au générateur d'adresses de
branchement.

La vérification fonctionnelle, conduite à l'aide de la suite
officielle \texttt{rv32ui-p}, montre que 40 des 42 tests passent
avec succès. Les deux échecs restants sont bien identifiés : le
test \texttt{ld\_st} (ID~39) met en évidence un cas limite dans
l'interaction entre le renvoi MEM-to-MEM et le mécanisme de
blocage, tandis que le test \texttt{ma\_data} (ID~41) confirme le
comportement attendu face à des accès mémoire mal alignés en
l'absence de gestionnaire de trappe. Ce taux de réussite de
95,2\,\% valide la conformité du c\oe{}ur pour l'intégralité du
sous-jeu utilisateur RV32I sur des programmes respectant
l'alignement naturel des accès mémoire.

Les perspectives identifiées sont, par ordre croissant
d'intrusion : l'ajout d'un prédicteur de branchement à 2~bits
indexé par PC ; l'implémentation d'un fichier CSR minimal et d'un
chemin de trappe pour la gestion des exceptions matérielles ;
l'intégration de l'extension \texttt{M} via un multiplieur et un
diviseur multi-cycles parallèles à l'ALU ; et enfin la
transformation des ports mémoire en ports AXI4-Lite, ouvrant la
voie à une intégration complète dans un système sur puce.
Chacune de ces extensions reste compatible avec la structure
existante du c\oe{}ur, ce qui souligne la modularité de la
conception présentée.

%==============================================================================
\appendix
\chapter{Liste des figures et descriptions}
\label{app:figures}

\noindent Les emplacements suivants correspondent, dans l'ordre, à
chaque balise de figure insérée dans le corps du
rapport. Chaque entrée est réservée pour une description technique
détaillée de la figure correspondante.

\begin{description}
  \item[Figure 1 --- Schéma bloc de plus haut niveau.]\hfill\\
        Entrées-sorties de plus haut niveau, interfaces IMEM/DMEM,
        chemin de l'adresse de démarrage, sortie
        \texttt{exception\_o}.

  \item[Figure 2 --- Hiérarchie des modules de
        \texttt{mpt\_top}.]\hfill\\
        Modules d'étage et blocs partagés (décodage, contrôle, banc
        de registres, ALU, génération de branchements).

  \item[Figure 3 --- Chemin de données pipeliné.]\hfill\\
        Chemin de données à cinq étages, registres pipeline
        inter-étages, multiplexeurs de renvoi.

  \item[Figure 4 --- Chemin de contrôle.]\hfill\\
        Génération du mot de contrôle à l'étage ID, propagation à
        travers les registres pipeline ID/EX, EX/MEM et MEM/WB.

  \item[Figure 5 --- Détail de l'étage IF.]\hfill\\
        Multiplexeur 4-vers-1 d'entrée du PC, duplication
        \texttt{IMEM\_addr\_o}/\texttt{IF\_pc\_o}, multiplexeur
        d'injection de NOP par \texttt{flush\_i}.

  \item[Figure 6 --- Architecture interne de l'étage ID.]\hfill\\
        Décodeur, banc de registres, toujours-blocs synchrones de
        capture, contrôle avec politique de bulle, multiplexeur de
        renvoi ID-WB.

  \item[Figure 7 --- Vue interne de
        \texttt{mpt\_branchgen}.]\hfill\\
        Multiplexeur \texttt{PC\_RELATIVE}/\texttt{REG\_OFFSET},
        additionneur 32~bits, renvoi local pour \texttt{JALR}.

  \item[Figure 8 --- Vue interne de l'étage MEM.]\hfill\\
        Décomposition adresse/donnée, génération de
        \texttt{DMEM\_wr\_byte\_en\_o}, unité d'extraction et
        d'extension, multiplexeur \texttt{ForwardM}.

  \item[Figure 9 --- Chronogramme d'aléa
        charge-utilisation.]\hfill\\
        Trois cycles consécutifs montrant le \texttt{LW}, la bulle
        insérée et le renvoi MEM/WB.

  \item[Figure 10 --- Vue synthétique de la stratégie
        d'aléas.]\hfill\\
        Sources de \texttt{forwardA}, \texttt{forwardB} et
        \texttt{forwardM}, chemin de \texttt{stall\_o}, fenêtre de
        \textit{flush} sur deux cycles.

  \item[Figure 11 --- Sortie du testbench en mode
        régression.]\hfill\\
        Capture d'écran des 42 verdicts \texttt{[PASS]} et du
        récapitulatif final.

  \item[Figure 12 --- Forme d'onde GTKWave d'un test
        typique.]\hfill\\
        Montée de \texttt{x3}, \texttt{x17} et \texttt{x10} aux
        valeurs de signature de succès.
\end{description}

\end{document}
